\section{Corrupted records}\label{sec:contamination}

\subsection{False alarms from a handful of records}\label{sec:c1}

Table~S5 of the Supporting Information gives the false-alarm rate of every test under corruption C1: $k$ records whose
covariate has been multiplied by four, in a sample for which the model is otherwise correct. Over the
$24$ scenarios of C1 the slope on the corrupted covariate moves by less than $5\%$, the slope on the
second covariate by at most $1.9\%$ and the intercept by at most $0.02$ on the linear-predictor scale,
so the model under test is effectively the clean one; under the sign error the slope moves by up to
$14\%$ (Section~\ref{sec:mechanism}). The fitted risks of the clean records do not move either. Their
mean absolute distance from the true risk, an integrated calibration index \citep{austin2019ici}, is
$0.0185$ with one exaggerated record in $1000$ and $0.0199$ with ten, against $0.0189$ without
corruption, which is the estimation error of any fit at this size; the corrupted records
themselves are mispredicted by about $20$ percentage points. Ten sign errors move it to $0.027$. A test
that rejects here is reacting to a few records, not to the model that is used for everyone else.

One corrupted record in a thousand moves every record-level test. Stukel's joint score rises from the
nominal level to $0.121$, its one-parameter form to $0.115$, the cubic calibration test to $0.127$ and
the GiViTI belt to $0.089$. EDGE reads $0.056$ with the cubic basis and $0.058$ with the
symmetric basis. At five corrupted records the record-level tests are between $0.25$ and $0.40$ and 
EDGE is still at $0.057$. At ten, Stukel's joint score reaches $0.627$, rejecting a correct
model more often than not, while EDGE reads $0.097$. Every test holds its level in the
same design without corruption (the $k=0$ rows), so the difference is not that one test is
conservative.


The pooled tests are protected, as Section~\ref{sec:mechanism} predicts. Le Cessie's smoothed test
reads $0.055$, $0.052$, $0.052$ and $0.077$ at one, two, five and ten corrupted records: an isolated
record adds one term to a sum whose spread is set by all $n$ records. The Hosmer--Lemeshow statistic and
the Pigeon--Heyse test at ten groups are protected too: under the exaggeration they stay at or below
$0.058$ up to ten corrupted records in $1000$ and twenty-five in $5000$ (Table~\ref{tab:severity}),
although the Pigeon--Heyse test is conservative throughout, between $0.019$ and $0.035$ on clean data.
Other grouped tests are not protected. The equal-width Hosmer--Lemeshow statistic, whose extreme groups hold few
records, reaches $0.592$ at ten corrupted records in $1000$, and the two tests that group on the
covariates, whose outlying clusters can be small, reach $0.542$ (Tsiatis) and $0.606$ (Xie). In a smaller study of three corrupted
records in $500$ with the covariate multiplied by four or by eight, $100$ replicates a scenario, BAGofT
and EDGE with the cubic basis gave no sign of moving, while the projection test reached
$0.22$ at $\times8$ and EDGE with the symmetric basis $0.20$; with $100$ replicates a rise
of a few points could not be seen (Supporting Information S3).

Wrong outcomes are a different matter. When the ten highest-risk records in a
thousand have their event erased (C2), every test rejects more often, EDGE included, at
$0.667$ against Stukel's $0.718$ and the cubic test's $0.837$. Missing events at the top of the risk
scale are a miscalibration of the recorded data, and a calibration test should find them. A unit error
that makes predictions milder (C3) moves no test under the logistic truth, although under the two
misfit truths it costs some power, and duplicated records (C4) leave every test between $0.040$ and
$0.071$ (Supporting Information S2).

\subsection{What protects: pooling with unit weights}\label{sec:whatprotects}

Proposition~\ref{prop:bounded} contrasts a grouped statistic with an ungrouped direction that grows with
the linear predictor. A direction that is bounded, a polynomial in the predicted risk rather than in the
logit, might protect without grouping, and the last study of Section~\ref{sec:contamination-design}
tests this on the same data sets (Supporting Information S10). It does not. With one exaggerated record
in $1000$ the ungrouped score test for the cubic directions of EDGE rejects a correct model in $0.166$
of replicates and with ten in $0.641$, against $0.047$ and $0.101$ for EDGE at the default partition and
$0.051$ and $0.066$ at ten groups; with ten reversed records it rejects in $0.992$. EDGE with one record a
group behaves like Stukel's test ($0.148$ and $0.594$). Spiegelhalter's $z$ never rejects at $5\%$ on
development data, because its normal reference ignores the fit, but read at its own realised level it
rejects in $0.118$ and $0.546$. Grouping alone is not enough either: in the study of
Section~\ref{sec:c1}, EDGE's score form, which uses the same groups but divides by the information its
directions keep after the fit, rejects in $0.087$ and $0.463$ at ten groups with one and ten
exaggerated records, against $0.056$ and $0.085$ for the unit form. The weighted grouped tests of \citet{hosmer2002goodness}, which weight each
record within its group by a function of its covariates, reach $0.18$ and $0.27$ with ten exaggerated
records and $0.68$ and $0.82$ with ten reversed ones. On frozen predictions, where nothing is refitted, the picture changes, which locates the damage
(Supporting Information S10). The same ungrouped twin then reads $0.061$ and $0.096$ with one and ten
exaggerated records in $1000$, and Spiegelhalter's $z$ holds throughout ($0.090$ with ten reversed
records), as in Section~\ref{sec:realcorrupt}. Two things therefore make a test fragile: the refit,
whose effect a score test magnifies by dividing by the information left after it, and, without a refit,
directions that lean on the few most extreme records, since the twin on frozen predictions still reads
$0.311$ with ten reversed records, where EDGE at ten groups reads $0.082$. EDGE's unit form avoids both:
its divisor is the variance of the clean groupmates, which neither the record nor the refit can make
small. The same holds with twelve
covariates at $n=2000$: with ten reversed records in one covariate EDGE reads $0.091$ at the default
partition and $0.048$ at ten groups, against $0.315$ for Stukel's test and $0.292$ for the ungrouped twin.

A routine alternative is to find the corrupted records with influence diagnostics and test on the rest.
Removing every record whose $\Delta X^2$ exceeds $4$ or whose $\Delta\beta$ exceeds $1$
\citep{pregibon1981logistic,hosmer2013applied} removed about $3.7\%$ of the records of a clean data set and
made every test reject a correct model in more than $0.95$ of replicates; removing only the records
with $\Delta\beta>1$ removed none, corrupted or not, and changed nothing. At an extreme prediction a
corrupted record's leverage and residual are both small on the fitted scale, so the diagnostics do not
single it out.

\subsection{The partition, not the contamination rate}\label{sec:partition}

On the same data the Hosmer--Lemeshow statistic reads $0.437$ at fifty corrupted records in $5000$ at
the default partition, where groups hold twenty-five records, and $0.088$ at ten groups, where they
hold five hundred; EDGE at the default partition reads $0.740$, and the record-level
tests between $0.87$ and $0.98$.

A second study varied the number of groups with the data held fixed, so that the same replicate is
tested at every partition (Figure~\ref{fig:partition}). The contamination \emph{rate} does not decide
the outcome: twenty-five corrupted records in $5000$ and five in $1000$ are both $0.5\%$, and 
EDGE reads $0.152$ and $0.048$, whereas the partition changes it: at $n=5000$ with twenty-five corrupted
records the same data give $0.062$ at ten groups and $0.152$ at two hundred; with fifty they give
$0.094$ and $0.726$. Without corruption the rate stays between $0.046$ and $0.057$ at every
partition.

\begin{figure}[htbp]\centering
\includegraphics[width=\textwidth]{Fig12_partition}
\caption{False-alarm rate of EDGE (cubic basis, unit form) under a correct model when $k$
records carry a corrupted covariate, against the number of risk groups. Each line is a fixed $k$; the
shaded area lies below $0.10$, where the test is read as holding, the dotted vertical line marks the
default partition, and the dashed grey line is the uncorrupted control. (a)~$n=1000$; (b)~$n=5000$. The
same contamination rate, $0.5\%$, gives different false-alarm rates in (a) and (b): the size of the
group, not the rate, decides.}
\label{fig:partition}
\end{figure}

The reason is in \eqref{eq:twomoments}. Corrupted covariates produce the most extreme predictions, so
they fill the top and bottom groups first; about $0.35k$ to $0.38k$ of them reach the top group at the
default partition, and about $0.40k$ at ten groups. The shift in that group's residual grows as
$\theta m_c/\sqrt{(m-m_c)\bar v}$, which rises steeply as $m_c$ approaches $m$. At the default partition
the group size $m$ stays near twenty-five, so the \emph{count} of corrupted records decides: in the
partition study ten in $1000$ put an average of $3.7$ into the top group and the test reads $0.079$,
and fifty in $5000$ put $17.8$ there, leaving about seven clean members, and it reads $0.726$. At ten
groups $m=n/10$ grows with the sample; the same data put $4.0$ and $19.9$ records into top groups of
$100$ and $500$, and the test reads $0.053$ and $0.094$. The partition study drew its own data sets, so
its rates differ from those of Table~S5 by Monte Carlo error.

\subsection{The severity of the error}\label{sec:severity}

A sign error (C1$'$) reverses the prediction instead of exaggerating it: a high-risk patient is placed
at the bottom of the risk scale with the outcome intact, and $\theta$ in \eqref{eq:twomoments} rises
from about $0.3$ to about $0.7$. The fit also moves more, by $14\%$ in the slope, and the two channels
were not separated. Table~\ref{tab:severity} gives both errors. With the cubic basis at
the default partition EDGE holds to about two corrupted records in $1000$ and five in
$5000$ under the sign error, against about ten under the exaggeration: at ten records in $1000$ it
reads $0.381$. Two settings do much better. With ten groups the cubic basis holds to five records in
$1000$ and ten in $5000$. The symmetric basis at the default partition holds to ten in both, reading
$0.098$ and $0.058$, although against many exaggerated records it is no better than the cubic basis
($0.745$ at fifty in $5000$). The symmetric basis at ten groups is the exception to the partition
lever: under the sign error it is less protected than at the default partition ($0.431$ against
$0.098$ at ten records in $1000$), which we observe but cannot explain from \eqref{eq:twomoments}.
Whatever the setting, the ordering of
the tests is unchanged. At one corrupted record in a thousand EDGE reads $0.054$ against
Stukel's $0.364$, the cubic test's $0.389$ and the belt's $0.268$; at five, $0.124$ against $0.89$,
$0.89$ and $0.80$.

\input{tab_severity}

\subsection{Protection when records may be wrong}\label{sec:edgefr}

For the cubic basis the partition sets a trade-off between power and protection, and it is chosen
before the data are seen. The default partition has more power: on clean data it leads ten groups by
$0.017$ on average over the $158$ scenarios ($0.016$ to $0.018$), and by much more when the misfit is
confined to the extreme risks (Supporting Information S10). With the true risk of the top $2\%$ of
patients $40\%$ below its prediction its size-adjusted power is $0.42$ at $n=1000$, against $0.19$ at ten
groups; with that risk $20\%$ below its prediction at $n=5000$, $0.62$ against $0.29$; and with the ten
highest-risk outcomes erased (C2), $0.60$ against $0.26$. Ten groups tolerate more corrupted records: they read $0.085$ with
ten exaggerated records in $1000$ and $0.107$ with fifty in $5000$, where the default partition reads
$0.097$ and $0.740$, and under the sign error $0.081$ with five records in $1000$ and $0.077$ with ten in
$5000$, where the default partition holds to about two in $1000$. The partition is therefore chosen in the
analysis plan, from what is known about how the data were collected, and not after the results are
seen: running both and reporting the more convenient would bring back the dependence on $G$ that
Section~\ref{sec:realpartition} shows for the Hosmer--Lemeshow test. The software reports both; if only
the default partition rejects, the signal lies in the extreme groups, and their records should be
checked.

At ten groups the plain Hosmer--Lemeshow statistic holds its level further than EDGE
once there are more than a few reversed records: $0.099$ against $0.182$ at ten in $1000$, and $0.126$
against $0.249$ at twenty-five in $5000$. Up to five reversed records in $1000$ both stay below $0.10$
($0.081$ and $0.056$), although the Hosmer--Lemeshow statistic already rejects less often on the same
data sets (exact McNemar $p=0.013$).
The Pigeon--Heyse test holds further still, but it is conservative on clean data. What EDGE
offers at ten groups is power, $0.411$ against $0.267$ (Section~\ref{sec:census}).

Pooling only the two extreme groups and keeping the default partition in the middle matched ten groups
up to about ten corrupted records in $1000$ and fifty in $5000$, was less protected beyond, and had the
same power; removing the extreme records instead removes the signal with them (Supporting Information
S5).

Section~\ref{sec:whentouse} turns these results into a rule.

\subsection{The cost in power, and robust alternatives}\label{sec:cost}\label{sec:robustest}

Under corruption a rejection can come from the misfit or from the corrupted records. The raw rejection
rate barely moves within the protected range: against a complementary log--log truth at $n=1000$ EDGE
rejects in $0.818$ of replicates on clean data, $0.785$ with one corrupted record and $0.802$ with ten.
Part of that is false alarm, so power is also read at the critical value of the test's own
\emph{contaminated} null, the logistic truth at the same sample size with the same records corrupted,
which the analyst cannot know but which isolates the misfit. Read this way EDGE falls from $0.826$
to $0.771$ at one corrupted record and $0.648$ at ten; against a probit truth at $n=5000$ it falls
from $0.278$ to $0.233$ at five and $0.207$ at ten. Corrupted records and a real tail misfit occupy the
same extreme groups, so the averaging that absorbs the first also weakens the second. Neither basis pays
consistently less: read at the contaminated null, at five corrupted records in $5000$ against the
probit truth the symmetric basis loses nothing and the cubic basis $0.045$, while against the
complementary log--log truth at one and two records in $1000$ the symmetric basis loses more ($0.084$
against $0.055$; $0.153$ against $0.096$). Beyond the protected range the test has
no power left: at twenty-five corrupted records in $5000$ its raw rejection rate against the probit
truth, $0.143$, is below its false-alarm rate on a correct model, $0.151$, and read at its
contaminated null it is $0.045$.

A natural objection is that corrupted records call for a robust fit. We fitted every contaminated data
set both by maximum likelihood and by the Huber-type quasi-likelihood estimator of
\citet{cantoni2001robust}, which bounds the Pearson residual (Huber's $\psi$ at $1.345$) but puts no
weight on the covariates, and applied every test to both fits. After a robust fit, five of ten tests no longer hold their level on clean data: both Stukel
score tests ($0.199$ and $0.223$) and three of EDGE's four forms, the symmetric basis in
its unit form ($0.164$) and both score forms. Of the five that do, none has its false alarms reduced: at
ten corrupted records the GiViTI belt reads $0.379$ after maximum likelihood and $0.400$ after the
robust fit, the cubic calibration test $0.465$ and $0.466$ (on this study's own data sets), and EDGE $0.089$ and $0.087$.
A corrupted record's prediction is formed at its corrupted covariate, so refitting does not reach it;
restoring the slope, if anything, makes that prediction more extreme. Refitting the model robustly is therefore not a substitute for a test
whose own statistic bounds a record's influence \citep{hampel1986robust,heritier1994robust}
(Supporting Information S4).

The fair robust competitor is therefore a robust \emph{test}. We compared EDGE with the
robust quasi-deviance test of \citet{cantoni2001robust} for Stukel's two constructed covariates, on
$1000$ fresh data sets a scenario at $n=1000$, in two forms: a Huber form, which bounds the Pearson
residual, and a Mallows form, which also down-weights records of high leverage and is the form meant
for wrong covariates. Both hold their level on clean data ($0.049$ and $0.047$). The Mallows form holds
below $0.10$ with one or two corrupted records, although with two reversed records it already rejects
more often than EDGE on the same data sets ($0.075$ against $0.056$, exact McNemar
$p=0.004$), and it does not hold beyond: with ten
exaggerated records it reads $0.168$ and with ten reversed ones $0.789$, where EDGE at ten
groups reads $0.063$ and $0.192$ on the same data sets (exact McNemar $p<10^{-20}$ in both). The Huber
form is less protected still ($0.205$ and $0.923$). On clean data the robust tests are also less
powerful: against a complementary log--log truth their size-adjusted power is $0.610$, against 
EDGE's $0.779$ at ten groups and $0.795$ at the default partition. Robustifying the
record-level test protects it against one or two exaggerated records and one reversed one, at a cost
in power; pooling protects further
at no such cost.
